\section{Referencial teórico e trabalhos relacionados}
\label{referencial}

\pendente{A equipe já escreveu as Seções 1 e 2 no Overleaf. Substitua ou complemente esta seção com o texto e as referências da revisão (T1 a T7). Abaixo ficam só os conceitos usados no método, com referências conferidas.}

\subsection{Detectores leves de objetos}

Detectores de estágio único fazem a localização e a classificação em uma só passagem pela rede, o que os torna adequados a dispositivos com pouca capacidade de processamento. O SSD \citep{liu2016ssd} prevê caixas a partir de âncoras distribuídas em mapas de características de várias escalas; combinado à MobileNet \citep{howard2017mobilenets}, que usa convoluções separáveis em profundidade para reduzir o custo, forma o SSD-MobileNet. O EfficientDet \citep{tan2020efficientdet} usa uma rede de fusão de características em várias escalas e escala largura, profundidade e resolução de forma conjunta; a variante Lite0 é a menor da família. O YOLOv8n é a menor variante da família YOLO distribuída pela Ultralytics \citep{jocher2023yolo}, sem âncoras e com resolução de entrada definida na exportação.

\subsection{Quantização INT8}

A quantização converte pesos e ativações de ponto flutuante de 32 bits para inteiros de 8 bits, com uma escala e um ponto zero por tensor, o que reduz o tamanho do modelo e permite aritmética inteira \citep{jacob2018quantization}. Na quantização pós-treinamento, os intervalos das ativações são estimados com um pequeno conjunto de calibração, sem novo treinamento. Os modelos deste trabalho rodam no LiteRT, o tempo de execução da Google para inferência no dispositivo \citep{litert2026}. Aceleradores como o Edge TPU exigem modelos inteiramente quantizados e só executam as operações que o compilador consegue mapear \citep{coral2026compiler}.

\subsection{Conjunto de dados e métrica de qualidade}

O COCO \citep{lin2014coco} é o conjunto de referência para detecção de objetos e inclui a classe pessoa. A métrica adotada é a precisão média (AP) com limiar de interseção sobre união de 0,5 (mAP@50), calculada com a ferramenta oficial do conjunto \citep{cocoapi2026}.

\subsection{Trabalhos relacionados}

\pendente{Colar o texto da equipe: busca sistemática (strings, bases, período, 3.362 registros), os 7 trabalhos-chave comparados e as lacunas L1 a L5. Pontos levantados nos slides de progresso, a referenciar: na CPU do Raspberry Pi 4, o SSD-MobileNet V1 leva cerca de 209 ms por imagem e o YOLOv8, de 760 ms a mais de 3 s; o Edge TPU reduz o SSD para cerca de 10 ms, mas obriga o YOLOv8n a usar entrada de 320 pixels (T1: Alqahtani, Cheema e Toosi, 2024); a quantização aproxima os modelos leves do tempo real com perda pequena de mAP (T2: Surantha e Sutisna, 2025). Incluir as referências completas no .bib.}
